% =====================================================================
%  c01 —— PDF p.9–16（原书 p.3–10）
%  §1 Stable curve 与 moduli stack；§1.1–§1.3
% =====================================================================

\section{Stable curve 的 moduli stack}
\subsection{Moduli space 与 moduli stack}

% ===== 原书 p.3 =====
本章讨论如何给代数曲线的同构类全体赋予代数簇乃至其推广的结构。先乐观地假定：给定亏格$g$的代数曲线全体（构成的空间）暂且记作$\mathcal{M}_g$，再假定它带有代数簇的结构，并且是万有的。要说明何谓万有，必须先定义亏格$g$代数曲线的「族」。代数簇的态射$f:C\to S$称为$S$上亏格$g$代数曲线的族，是指$f$为紧合光滑态射，且各纤维都是连通的一维代数簇，其亏格均为$g$。此时称$S$为该族的底簇（base variety），称$f$为投影（projection）。$\mathcal{M}_g$是万有的，是指：

\begin{quote}
存在$\mathcal{M}_g$上亏格$g$的代数曲线族$C_g\to\mathcal{M}_g$，使得任意底簇$S$上亏格$g$的代数曲线族$C\to S$，都可作为$C_g\to\mathcal{M}_g$沿某个态射$\phi:S\to\mathcal{M}_g$的拉回得到；并且具有这一性质的$\phi$是唯一的。
\end{quote}

在这样的框架中，基础域的特征不必为$0$；但为简单起见，以下取基础域为$\C$。若这样的族$C_g\to\mathcal{M}_g$存在，那就很幸运了。因为由一点到$\mathcal{M}_g$的态射全体与亏格$g$代数曲线的同构类全体一致；更一般地，代数簇$S$上亏格$g$代数曲线的同构类全体与$S$到$\mathcal{M}_g$的态射全体一致。特别地，若另有具此性质的$C'_g\to M'_g$，则它与$C_g\to M_g$通过唯一的同构等同起来。然而一般地，代数曲线具有自同构，因此不能指望这样的族存在。例如设$C$为具有$d$阶自同构$\sigma$的代数曲线。此时令$\zeta=\exp(2\pi i/d)$，用$\tilde\sigma=\sigma\times\zeta:(x,t)\mapsto(\sigma(x),\zeta t)$作用在$C\times\C^\times$上，就得到$\Z/d\Z$的作用。将此作用在$C\times\C^\times$上的商记作$\tilde C$。自然投影给出代数曲线族$\tilde C\to\C^\times/\langle\zeta\rangle$；其各纤维都与$C$同构，但它不同构于直积。（称此为$\sigma$扭过的族。）若存在万有族$C_g\to\mathcal{M}_g$，则$\C^\times/\langle\zeta\rangle$到$\mathcal{M}_g$的态射必须是常值映射。把该常值映射拉回，若所得族与$\tilde C\to\C^\times/\langle\zeta\rangle$同构，那就成了直积族，矛盾。自同构的存在构成了障碍；但不幸中的万幸是，亏格不小于$2$的曲线，其自同构群总是有限群。于是可以期望：给曲线加上标记，使它不至于「乱转」，用这个办法来回避困难。带着标记，却仿佛并不意识到标记；或者说与标记无关的表述，所用的就是代数叠（algebraic stack）。更准确的定义留待后面，先讲曲线退化与模空间紧化所不可缺少的稳定曲线的定义与性质。

\subsection{关于 stable curve}

% ===== 原书 p.4 =====
本节讲述 Deligne--Mumford [DM] 关于曲线模空间紧化的结果。设$g$为不小于$2$的自然数。

\begin{definition}
$f:C\to S$称为亏格$g$的稳定曲线（stable curve），是指它满足以下条件：
\begin{enumerate}[label=(\arabic*)]
\item $f$紧合平坦，且相对维数为$1$。
\item 各几何点$s$处的几何纤维$C_s$既约连通，且至多只有通常二重点。
\item 每个非奇异有理曲线分支$R$与其它分支的交点至少有$3$个。
\item $\dim_{k(s)}H^1(C_s,\mathcal{O}_{C_s})=g$。
\end{enumerate}
\end{definition}

\begin{theorem}
设$k$为域，$C\to\Spec(k)$为$\Spec(k)$上的稳定曲线。令$\omega_{C/k}$为对偶层（dualizing sheaf），则$\omega_{C/k}^{\otimes3}$极丰沛。
\end{theorem}

\begin{proof}
设$\pi:\tilde C\to C$为正规化，$\Sigma$为$C$的奇点全体，令$\tilde\Sigma=\pi^{-1}(\Sigma)$。要说明$\omega^{\otimes3}$极丰沛，只需对$x,y\in C$说明
\[
(*)\qquad H^0(C,\omega_C^{\otimes3})\longrightarrow H^0(C,\omega_C^{\otimes3}\otimes\mathcal{O}_C/\mathcal{M}_x\mathcal{M}_y)
\]
为满射即可。为简单起见，考虑$x,y\in C-\Sigma$且存在$\tilde C$的分支$E$使$x,y\in\pi(E)$的情形。此时只需说明
\[
\dim H^0(C,\omega_C^{\otimes3}(-x-y))=\dim H^0(C,\omega_C^{\otimes3})-2
\]

% ===== 原书 p.5 =====
即可。考虑正合列
\[
0\to\omega_C\to\pi_*(\Omega^1_{\tilde C}(\tilde\Sigma))\to k(\Sigma)\to0,
\]
便得到
\begin{align*}
H^0(C,\omega_C^{\otimes3})&=\ker\bigl(H^0((\Omega^1_{\tilde C}(\tilde\Sigma))^{\otimes3})\to k(\Sigma)\bigr),\\
H^0(C,\omega_C^{\otimes3}(-x-y))&=\ker\bigl((H^0((\Omega^1_{\tilde C}(\tilde\Sigma))^{\otimes3}(-x-y))\to k(\Sigma)\bigr).
\end{align*}
因此令$\Sigma_E=\tilde\Sigma\cap E$后，只需看
\[
H^0(E,(\Omega^1_E)^{\otimes3}(3\Sigma_E))\longrightarrow(\mathcal{O}_E/\mathcal{M}_x\mathcal{M}_y)\oplus k(\Sigma_E)
\]
为满射。为此用 Serre 对偶，只需说明
\[
(**)\qquad H^0(E,(\Omega^1_E)^{\otimes(-2)}(-2\Sigma_E+x+y))=0
\]
即可。由$\deg(\Omega^1_E(\Sigma_E))\ge1$，只要$\deg(\Omega^1_E(\Sigma_E))\ne1$，$(**)$中线丛的次数就是负的，故$(**)$成立。若$\deg(\Omega^1_E(\Sigma_E))=1$，则$(**)$中线丛的次数为$0$；考察它与其它分支的交，那里必须消失，所以$(**)$仍然消失。（Deligne--Mumford 在论文[DM]中漏掉了$C=E_1\cup E_2$，$p(E_i)=1$（$i=1,2$），$E_1\cap E_2=\{p\}$，$x=y=p$的情形，但此时同样可以说$(*)$是满射。）
\end{proof}

\transnote{原文此处写作「$p(E_i)=1$」，其中$p$的用法在书中未加说明，疑为亏格记法。译文照原文保留。}

现在令$V=H^0(C,\omega_{C/k}^{\otimes3})$，$\mathbf{P}=\mathbf{P}(V)$。此时三次典范映射（tricanonical map）$\Phi:C\to\mathbf{P}$是嵌入，故可把$C$看作$\mathbf{P}$的闭子簇。它的 Hilbert 多项式$f(m)=\chi(C,\omega_{C/k}^{\otimes3m})=\dim H^0(C,\omega_{C/k}^{\otimes3m})$由
\[
p(m)=3m(2g-2)+1-g=(6m-1)(g-1)
\]
给出。这里先复习一下 Hilbert 概形。

% ===== 原书 p.6 =====
设$W$为$n+1$维向量空间。对$\mathbf{P}=\mathbf{P}(W)$的子概形$V$，定义其 Hilbert 多项式为
\[
p(m)=\chi(\mathbf{P},\mathcal{O}_V(m))=\sum_i(-1)^iH^i(\mathbf{P},\mathcal{O}_V(m)),
\]
则它是次数为$\dim V$的多项式。此外，已知以下事实。

\begin{enumerate}[label=(\arabic*)]
\item 设$V$为$\mathbf{P}\times S$的闭子簇且在$S$上平坦。此时在$S$的连通分支上，$S$的点$s$的各纤维$V_s$的 Hilbert 多项式$p(m)$与$s$无关。称这样的$V$为$S$上以$p(m)$为 Hilbert 多项式的平坦族。
\item 设$p(m)$为给定的关于$m$的多项式。此时存在只依赖于$p$的$m_0$，使得对大于$m_0$的$m$，下述事实成立。对任意以$p$为 Hilbert 多项式的$S$上平坦族$V$，
\[
p_2^*p_{2*}\mathcal{I}(m)\longrightarrow\mathcal{I}(m)
\]
为满射，且$R^ipr_{2*}\mathcal{I}(m)=0$（$i>0$）。这里$pr_2$为自然投影$\mathbf{P}\times S\to S$，$\mathcal{I}$为$V$的理想层。
\end{enumerate}

\transnote{原文公式中写作 $p_2^*p_{2*}\mathcal I(m)$，同段正文却用 $\mathrm{pr}_2$ 表示自然投影；公式与正文记号混用，二者指同一态射。译文照原文保留。}

\begin{theorem}[Grothendieck]
设$p$为给定的多项式。对$k$上的概形$S$，对应集合
\begin{align*}
\mathrm{Hilb}^p(S)=\{V\subset\mathbf{P}\times S\mid{}&V\text{ 为以 }p\text{ 为 Hilbert 多项式的}\\
& S\text{ 上平坦族}\}
\end{align*}
的函子可表。即存在概形$\mathrm{Hilb}^p$以及$\mathbf{P}\times\mathrm{Hilb}^p$的万有闭子簇$V^p$，使得$\mathrm{Hilb}^p(S)$的任意元素$V\to S$，都可由唯一的态射$S\to\mathrm{Hilb}^p$把$V^p\to\mathrm{Hilb}^p$拉回得到。
\end{theorem}

\begin{proof}[证明概略]
对$p$取上述事实(2)中确定的$m_0$，令$a=p(m_0)$。设$\mathrm{Sym}_{m_0}(W)$为$W$的$m_0$阶对称积，$G$为$\mathrm{Sym}_{m_0}(W)$的$a$维子空间构成的 Grassmann 簇。$G$上存在余秩$a$的万有向量子丛$\mathcal{U}\subset\mathcal{O}_G\otimes\mathrm{Sym}_{m_0}(W)$；把$G\times\mathbf{P}$上由$\mathcal{U}$生成的理想层所对应的闭子概形记作$\tilde V$。对于$\tilde V\to G$，存在称为平坦化分层（flattening stratification）的下述分层：$G=\bigsqcup_jG^{p_j}$，其中$G^{p_j}$为局部闭子簇，且$\tilde V$到$G^{p_j}$的限制是以$p_j$为 Hilbert 多项式的平坦族。这里把与$p_j=p(m)$对应的$G^{p_j}$记作$\mathrm{Hilb}^p$，它就是所求的对象。此处须注意一点：函子$\mathrm{Hilb}^p$被全空间$\mathbf{P}$打上了标记（marking）。
\end{proof}

% ===== 原书 p.7 =====
现在回到稳定曲线。作为$\mathrm{Hilb}^p$的子函子，如下定义三次典范稳定曲线函子$H_g$：
\[
H_g(S)=\{(f:C\to S,\ \phi:\mathbf{P}(f_*\omega_{C/S}^{\otimes3})\cong\mathbf{P}^{5g-6})\mid f:C\to S\text{ 是稳定曲线}\}.
\]

\begin{theorem}[Deligne--Mumford {[DM]}]
$H_g$由$\mathrm{Hilb}^p$的子概形表示。令$Z_g$为万有曲线，则$H_g$与$Z_g$都由$k$上的光滑代数簇表示。
\end{theorem}

\begin{proof}
证明用 Schlessinger 的变形理论。取$H_g$的闭点$S$，由$\Ext^2(\omega_{C/S},\mathcal{O}_C)=0$可知，万有形式变形（versal formal deformation）是形式光滑的。再由$\Hom(\omega_{C/S},\mathcal{O}_S)=0$可知，万有形式变形是万有的。现取$\mathcal{X}\to\mathcal{M}=\Spec(k[[x_1,\ldots,x_m]])$为万有族，并取一个同构$H^0(\mathcal{X},\omega_{\mathcal{X}/\mathcal{M}})\cong\mathcal{O}_{\mathcal{M}}^{5g-5}$，便定出态射$\mathcal{M}\to H_g$；它经由$S$的形式邻域$T$分解。又由万有性，从$T$可定出$\mathcal{M}$；而$T$是以$\mathcal{M}$的像为一个截面的形式群概形$P\widehat{GL}(5g-5)$为结构群的形式主丛。因此$T$光滑。
\end{proof}

\transnote{原文此处把结构群写作「$P\widehat{GL}(5g-5)$」，帽子只覆盖$GL$，译文照原文保留。}

\subsection{关于 algebraic stack}

\begin{definition}
设$\mathcal{C}$为具有纤维积的范畴。$\mathcal{C}$上的 Grothendieck 拓扑，是指对$\mathcal{C}$的对象$X$给定了覆盖族（covering family）$\{U_i\to X\}_i$这一概念，并且具有以下性质：
\begin{enumerate}[label=(\arabic*)]
\item $\{X\to X\}$是$X$的覆盖族。
\item 若$\{U_i\to X\}_i$是$X$的覆盖族，且$\{U_{i,j}\to U_i\}_j$是$U_i$的覆盖，则其合成$\{U_{i,j}\to X\}_{i,j}$是$X$的覆盖族。
\item 若$\{U_i\to X\}_i$与$\{V_j\to X\}_j$都是$X$的覆盖族，则$\{U_i\times_XV_j\to X\}_{i,j}$是$X$的覆盖族。
\end{enumerate}
\end{definition}

Grothendieck 拓扑的例子有：拓扑空间中开覆盖构成的范畴，以及概形中平展满射族构成的范畴等。

\transnote{原文此处把 topological 拼作「topologiacal」，译文按「拓扑空间」译出。}

% ===== 原书 p.8 =====
又，映入具有 Grothendieck 拓扑的范畴$\mathcal{C}$的函子$p:\mathcal{S}\to\mathcal{C}$称为群胚纤维化（fibered in groupoid），是指它满足以下条件：
\begin{enumerate}[label=(\arabic*)]
\item 对$\mathcal{C}$的对象$V$，$\mathcal{S}$的子范畴
\[
\mathcal{S}_V=\{\text{对象}=X,\ \text{态射}=f\mid p(X)=V,\ p(f)=\id_V\}
\]
中所有态射都是同构。
\item 设$U,V$为$\mathcal{C}$的对象，$\phi:U\to V$为态射。又设$y$为$\mathcal{S}_V$的对象。此时存在$\mathcal{S}$的对象$x$与$\mathcal{S}$的态射$f:x\to y$，使得$p(f)=\phi$，$p(x)=U$。
\item 给定$\mathcal{S}$的对象$x,y,z$及态射$f:x\to z$，$g:y\to z$。对任意满足$p(g)\chi=p(f)$的$\chi:p(x)\to p(y)$，都唯一地确定$h:x\to y$，使得$p(h)=\chi$。
\end{enumerate}

\medskip
\noindent\textbf{补充.}
\begin{enumerate}[label=(\arabic*)]
\item 由(2)定出的$x$记作$\phi^*(y)$。
\item 至此尚未用到$\mathcal{C}$的 Grothendieck 拓扑。
\end{enumerate}

\begin{example}
考虑由下述对象与态射构成的范畴$\overline{\mathcal{M}}_g$：对象为亏格$g$的稳定曲线$f:C\to S$，其中$C$与$S$都是基础域$k$上有限型的；态射为交换图
\[
\begin{tikzcd}
C \arrow[r] \arrow[d, "f"'] & C' \arrow[d, "f'"] \\
S \arrow[r] & S'
\end{tikzcd}
\]
又令$\Sch(k)$为$k$上有限型概形构成的范畴。令函子$p$把$f:C\to S$对应到底概形$S$。此时$p:\overline{\mathcal{M}}_g\to\Sch(k)$是群胚纤维化。
\end{example}

\begin{example}
设$X$为$k$上有限型的概形。令$\Sch/X$为$X$上有限型概形构成的范畴。对$\Sch/X$的对象$Y\to X$，令$p$为把$Y$对应过去的函子。则$p:\Sch/X\to\Sch(k)$是群胚纤维化。此例称为可表叠（representable stack），特别重要。
\end{example}

\transnote{原文此处作「$X$ を $k$ 上有限型の Scheme する」，助词「と」下脱，按文意应为「…とする」。译文按「设为」译出。}

下面，$\mathcal{C}$取$\Sch(k)$并赋予适当的 Grothendieck 拓扑时特别重要。粗略地说，满足粘合条件的称为叠（stack）；进而局部平展地

% ===== 原书 p.9 =====
能像例 2 那样表示的，即可表的叠，称为代数叠（algebraic stack）。以下说明这一点。在此之前注意：$\mathcal{C}=\Sch(k)$上的群胚纤维化$p:\mathcal{S}\to\Sch(k)$也构成范畴（态射是什么此处不特别说明），而其中由可表叠构成的全子范畴自然等价于$k$上有限型概形的范畴。

\transnote{原文此处引作「Example 2」，指上文编号为 2 的可表叠之例；译文按全书例号体例改作「例 2」。}

\begin{definition}[stack]
设$\mathcal{C}$为具有 Grothendieck 拓扑（暂记作$e$）的范畴，$p:\mathcal{S}\to\mathcal{C}$为群胚纤维化。$p$为叠（stack），是指：
\begin{enumerate}[label=(\arabic*)]
\item 设$U$为$\mathcal{C}$的对象，$x,y$为$\mathcal{S}_U$的对象。对$(\mathcal{C}/U)_e$的对象$\phi:V\to U$，把$\Hom_{\mathcal{S}_V}(\phi^*(x),\phi^*(y))$对应过去的函子是$e$-拓扑下的层。

这里用了几个尚未定义的概念，为慎重起见先确认如下。$\mathcal{C}/U$是以$U$上的对象$\phi:V\to U$为对象的范畴；再对此范畴赋予与$\mathcal{C}$相同的覆盖族与 Grothendieck 拓扑（以下同一范畴上可能赋予两种以上拓扑，故此拓扑称为$e$-拓扑；$e$不必是 etale，但 etale 拓扑是典型的 Grothendieck 拓扑，故常用字母$e$）。$(\mathcal{C}/U)_e$上取值于集合范畴的反变函子$F$是（关于$e$-拓扑的）层，是指对任意对象$V\in(\mathcal{C}/U)$与任意覆盖族$h_i:V_i\to V$，成立下述两条：(1)把$F(V)$的元素$x$对应到$(h_i^*(x))_i\in\prod_iF(V_i)$的映射是单射；(2)若$\prod_iF(V_i)$的元素$(x_i)_i$满足$h_{i,j}^*(x_i)=h_{j,i}^*(x_j)$，则它属于(1)所给映射的像。这里$h_{i,j}$是自然映射$h_{i,j}:V_i\times_VV_j\to V_i$。
\item 设$h_i:U_i\to U$为$\mathcal{C}$中的覆盖族。若存在$\mathcal{S}_{U_i}$的对象$x_i$，且给定了$\mathcal{S}_{U_i\times_UU_j}$中满足上闭链条件（cocycle condition）的同构$\sigma_{i,j}:h_{i,j}^*(x_i)\cong h_{j,i}^*(x_j)$，则存在$\mathcal{S}_U$的对象$x$，使得$h_i^*(x)=x_i$。（称为有效下降条件（effective descent condition）。）
\end{enumerate}
\end{definition}

\medskip
\noindent\textbf{注意.}
$\mathcal{C}$上的叠具有纤维积。设$p:\mathcal{S}\to\mathcal{C}$，$q:\mathcal{T}\to\mathcal{C}$为两个叠。此时以
\[
\{(x,y,\phi)\mid x\in\mathcal{S},\ y\in\mathcal{T},\ \phi:p(x)\cong q(y)\}
\]
为对象的范畴，就是$\mathcal{S}$与$\mathcal{T}$在$\mathcal{C}$上的纤维积。

% ===== 原书 p.10 =====
以下设$\mathcal{C}$为$k$上有限型概形赋以平展拓扑所得的范畴。绕了一大圈，现在来定义代数叠。把概念构成所需的结构归纳起来，便得到下表。
\begin{itemize}[leftmargin=6em,itemsep=1pt,topsep=3pt]
\item 群胚纤维化 —— 范畴论
\item 叠 —— 带 Grothendieck 拓扑的范畴
\item 代数叠 —— 带平展拓扑的代数簇
\end{itemize}

\begin{definition}[algebraic stack]
设$S\to\Sch(k)$为$\Sch(k)$上的叠。$S$为代数叠（algebraic stack），是指：
\begin{enumerate}[label=(\arabic*)]
\item $\Delta:S\to S\times_{\Sch(k)}S$可表。即对任意可表叠$X\to\Sch(k)$以及$X$到$S\times_{\Sch(k)}S$的态射$X\to S\times_{\Sch(k)}S$，如下定义的纤维积$Z$可表：
\[
\begin{tikzcd}
Z \arrow[r] \arrow[d] & X \arrow[d] \\
S \arrow[r] & S\times_{\Sch(k)}S
\end{tikzcd}
\]
\item 存在可表叠$X$与平展局部满射的态射$X\to S$。即对$S$的任意对象$y$，存在$p(y)$的平展覆盖$h_i:U_i\to p(y)$，使得$h_i^*(y)$是$X_{U_i}$的像。
\end{enumerate}
\end{definition}

把上述两个条件就$\overline{\mathcal{M}}_g$来考虑，便得到如下情形。首先，当$X$是可表叠且由概形$X$表示时，给出函子$f:X\to\overline{\mathcal{M}}_g$与给出稳定曲线$C\to X$等价；这是因为$f$由$f(X\xrightarrow{\id}X)$唯一确定。

\begin{enumerate}[label=(\arabic*)]
\item 由$X$给出$\overline{\mathcal{M}}_g\times_{\Sch(k)}\overline{\mathcal{M}}_g$，等价于给出两条稳定曲线$C_1\to X$，$C_2\to X$。此时(1)中的范畴$Z$等价于四元组$((U\to X),(C\to U),\phi_1,\phi_2)$构成的范畴。这里$U\to X$是$X$上的概形，$C\to U$是$U$上的稳定曲线，$\phi_1,\phi_2$分别是
\[
(C\to U)\xrightarrow{\phi_1}(C_1\times_XU\to U),\qquad
(C\to U)\xrightarrow{\phi_2}(C_2\times_XU\to U)
\]
的同构。它由概形$\mathrm{Isom}_X(C_1,C_2)$表示。所谓概形$\mathrm{Isom}_X(C_1,C_2)$，粗略地说，是$X$上的概形，使各几何点$x\in X$处的纤维$\mathrm{Isom}_X(C_1,C_2)(k(x))$由$\mathrm{Isom}_{k(x)}(C_{1x},C_{2x})$给出。（严格地说，按惯例把它定义为表示某个函子的对象。）

% ===== 原书 p.11（承接上页段落，由本块译完）=====
\item 存在平展局部满射$X\to\overline{\mathcal{M}}_g$，等价于：存在$X$上的稳定曲线$C\to X$，使得对任意稳定曲线$D\to W$，都存在$W$的平展覆盖$W_i\to W$与态射$W_i\to X$，且对每个$i$，$D\times_WW_i\to W_i$与$C\times_XW_i\to W_i$同构。
\end{enumerate}

这里，由(2)的条件，代数叠不会自由得漫无边际。
